\title{xLSTM: Extended Long Short-Term Memory}

\begin{abstract}
During the 1990s, the pivotal concepts of the constant error carousel and gating mechanisms were established as foundational components of Long Short-Term Memory (LSTM) networks. LSTMs have since proven robust over time, playing essential roles in various deep learning advancements, most notably as the underlying architecture for the inaugural Large Language Models (LLMs). However, the emergence of the Transformer model—distinguished by its parallelizable self-attention mechanism—initiated a transformative period in sequence modeling and rapidly surpassed LSTM-based models in large-scale applications. In light of this, we pose a straightforward question: What performance is attainable in language modeling if LSTMs are scaled to billions of parameters while adopting the latest strategies from contemporary LLM research and addressing longstanding LSTM constraints? We address this by proposing exponential gating in conjunction with suitable normalization and stabilization methods. Furthermore, we alter the LSTM's memory architecture, leading to the development of: (i) sLSTM, which employs scalar memory and update mechanisms together with a redefined mixing operation, and (ii) mLSTM, characterized by complete parallelizability via a matrix memory and a covariance-based update mechanism. These advanced LSTM designs are embedded within residual block frameworks, resulting in xLSTM modules that can be composited into full xLSTM networks. Through exponential gating and adapted memory configurations, xLSTM architectures demonstrate superior performance and scalability in comparison to leading Transformers and State Space Models.\\
Code available at: \url{https://github.com/NX-AI/xlstm}
\end{abstract}

\section{Introduction}
\label{sec:intro}

The Long Short-Term Memory (LSTM) ideas~\citep{Hochreiter:91,Hochreiter:97nips,Hochreiter:97}, i.e., the constant error carousel and gating, were introduced to overcome the vanishing gradient problem of recurrent neural networks~\citep{Hochreiter:91,Hochreiter:00book}:

\begin{align}
\label{eq:lstm_idea}
  c_t \ = \ f_t \ c_{t-1} \ + \ 
          i_t \  z_t \ , \quad  
          h_t \ = \ o_t \ \psi({c_t}) \ . 
\end{align}

The constant error carousel functions by incrementally updating the cell state $c_{t-1}$ (shown in green) with new cell inputs $z_t$, while this process is regulated through sigmoid gating mechanisms (shown in blue). The update is modulated by both the input gate $i_t$ and the forget gate $f_t$, whereas the memory cell's output, or hidden state $h_t$, is governed by the output gate $o_t$. Prior to determining the hidden state, the cell state undergoes normalization or squashing via $\psi$, followed by application of output gating.

LSTMs have seen widespread adoption across multiple fields \citep{Hochreiter:01,Hochreiter:07,Schmidhuber:15}, maintaining dominance in text generation prior to the introduction of Transformers in 2017~\citep{Vaswani:17}. Their utility has been substantiated in a diverse set of sequential modeling tasks, including but not limited to text synthesis~\citep{Graves:13,Karpathy:15blog}, handwriting production~\citep{Graves:13}, sequence-to-sequence translation~\citep{Sutskever:14nips}, program evaluation~\citep{Zaremba:14arxiv}, automated image captioning~\citep{Karpathy:15,Hossain:19}, code generation~\citep{Karpathy:15blog}, hydrological modeling for rainfall-runoff prediction~\citep{Kratzert:18,Kratzert:19}, and operational flood warning systems~\citep{Nearing:24}. Within reinforcement learning, LSTMs continue to outperform other sequence modeling techniques, exemplified in prominent systems like AlphaStar for StarCraft II~\citep{Vinyals:17short}, OpenAI Five for Dota 2~\citep{OpenAI:19}, and controllers for magnetic confinement in nuclear fusion~\citep{Degrave:22short}. LSTMs possess a notable ability to learn complex abstractions, efficiently capturing and retaining semantic content via their memory mechanisms~\citep{Karpathy:15blog}. This capacity is highlighted by findings such as number and syntax neurons~\citep{Lakretz:19}, neurons that represent linguistic properties~\citep{Bau:19}, and neurons associated with sentiment~\citep{Radford:17arxiv}. Despite advances in sequence modeling, LSTMs persist in cutting-edge practical domains~\citep{Degrave:22short,Nearing:24}, demonstrating remarkable longevity and resilience.

Although LSTMs have achieved significant breakthroughs, they exhibit three principal shortcomings: (i) Inability to update storage choices post hoc. This constraint is illustrated using the {\it Nearest Neighbor Search} problem (refer also to Appendix~\ref{sec:appNearestNeighborSearch}): When provided with a reference vector, a sequential scan is required to locate the most similar vector and subsequently retrieve its associated value at the end of the sequence. As depicted in the left panel of Figure~\ref{fig:lstmProblems}, the task’s mean squared error demonstrates that LSTM architectures struggle to overwrite a previously stored value when a vector with greater similarity appears later in the sequence. Our xLSTM model addresses this with the use of exponential gating mechanisms. (ii) Constrained capacity for information retention, as data must be encoded within scalar-valued cell states. This restriction is evident in the context of {\it Rare Token Prediction}, depicted in the right panel of Figure~\ref{fig:lstmProblems}, where perplexity scores for Wikitext-103~\citep{Merity:17} are broken down by frequency bands. LSTMs underperform on infrequent tokens, largely due to their storage bottleneck. The introduction of a matrix-based memory in xLSTM successfully overcomes this issue. (iii) An absence of parallelization capabilities stemming from memory interdependence, specifically the hidden-to-hidden transitions between successive time steps, which enforces strictly sequential computation.

These limitations have catalyzed the rise of Transformer models~\citep{Vaswani:17} in the domain of language modeling. The question arises: What language modeling performance can be attained if these limitations are resolved and LSTMs are scaled to match the parameter size of current Large Language Models?

\section{Extended Long Short-Term Memory}
\label{sec:xLSTM}

To address the shortcomings of the standard LSTM, Extended Long Short-Term Memory (xLSTM) brings forward two significant changes to the core LSTM framework described in Equation~\eqref{eq:lstm_idea}. The two primary innovations—exponential gating and alternative memory architectures—give rise to two novel variants within the LSTM family: (i) the sLSTM (Section~\ref{sec:sLSTM}), which utilizes a single scalar memory, a scalar update, and integrates a process of memory mixing, and (ii) the mLSTM (Section~\ref{sec:mLSTM}), which employs a matrix memory and updates via an outer product (covariance-based) mechanism that supports full parallelization. Both sLSTM and mLSTM advance beyond traditional LSTM designs by incorporating exponential gating. In pursuit of improved parallelism, mLSTM dispenses with memory mixing, specifically removing the recurrent hidden-to-hidden connections. Both architectures, mLSTM and sLSTM, are capable of being generalized to encompass multiple memory cells. For sLSTM, this multi-cell approach enables memory mixing across different cells. Moreover, the sLSTM model can support multiple attention heads, wherein mixing is restricted to cells within a head and does not occur between heads. The addition of heads, alongside exponential gating, enables a novel scheme for memory mixing in sLSTM. Conversely, for the mLSTM architecture, having multiple heads or multiple cells is functionally equivalent.

By embedding these enhanced LSTM variants within residual block structures, one obtains xLSTM blocks (refer to Section~\ref{sec:xLSTMblock}). Building neural networks by stacking xLSTM blocks in a residual manner results in complete xLSTM architectures (see Section~\ref{sec:xLSTMarchitecure}). The overall layout of the xLSTM framework and its constituents is depicted in Figure~\ref{fig:xlstm_sketch}.

\subsection{Review of the Long Short-Term Memory}
\label{sec:orgLSTM}

The foundational concept of LSTM~\citep{Hochreiter:91,Hochreiter:97nips,Hochreiter:97} centers on a scalar memory cell, designed as the core unit for processing and retaining information while mitigating the vanishing gradient problem~\citep{Hochreiter:91,Hochreiter:00book} via the constant error carousel mechanism inherent in the cell state dynamics. This memory cell architecture incorporates three primary gates: the input gate, output gate, and forget gate, the latter of which was subsequently added by \citet{Gers:00}. The equations governing the LSTM memory cell update at each discrete time $t$ are:

\begin{align}
c_t \ &= \  f_t \ c_{t-1} \ + \ i_t \ z_t &  & &\text{cell state} \\
h_t  \ &= \ o_t \ \tilde{h}_t \ , & \tilde{h}_t \ &= \ \psi \left( c_t\right)
  &\text{hidden state} \\
z_t \ &= \ \varphi \left( \tilde{z}_t \right) \ , 
  &\tilde{z}_t \ &=  \ \mathbf{w}^\top_{z} \ \mathbf{x}_t \ + \
  r_{z}  h_{t-1} \ + \  b_{z} \ \
  &\text{cell input} \\
i_t \ &= \ \sigma \left( \tilde{i}_t  \right) \ , 
  &\tilde{i}_t \ &= \ \mathbf{w}^\top_{i} \ \mathbf{x}_t \ + \
  r_{i}  \ h_{t-1} \ + \  b_{i} \ \
  &\text{input gate} \\
f_t \ &= \ \sigma \left(  \tilde{f}_t \right) \ , 
  &\tilde{f}_t \ &= \ \mathbf{w}^\top_{f} \ \mathbf{x}_t  \ + \
  r_{f}  \ h_{t-1} \ + \  b_{f} \ \
  &\text{forget gate} \\
o_t \ &= \ \sigma \left( \tilde{o}_t \right) \ , 
  &\tilde{o}_t  \ &= \ \mathbf{w}^\top_{o} \ \mathbf{x}_t \ + \
  r_{o}  \ h_{t-1} \ + \  b_{o} \ \
  &\text{output gate} 
\end{align}

The vectors $\mathbf{w}_{z}$, $\mathbf{w}_{i}$, $\mathbf{w}_{f}$, and $\mathbf{w}_{o}$ represent the input weight parameters linking the input $\mathbf{x}_t$ to the cell input, input gate, forget gate, and output gate, respectively. The weights $r_{z}$, $r_{i}$, $r_{f}$, and $r_{o}$ denote the recurrent connections from the previous hidden state $h_{t-1}$ to the cell input, input gate, forget gate, and output gate. The biases $b_{z}$, $b_{i}$, $b_{f}$, and $b_{o}$ are each associated with their respective components. The activation functions $\varphi$ and $\psi$ typically employ $\tanh$ for the cell input and hidden state. The function $\psi$ serves to restrict the range of the cell state, thus preventing it from becoming unbounded. All gate activations utilize the sigmoid function $\sigma \left( x \right)= 1/(1+\exp(-x))$. In subsequent LSTM architectures, the aggregation of multiple one-dimensional memory cells into a vector $\mathbf{c}_t \in \mathbb{R}^{d}$ permitted the application of recurrent weight matrices $\mathbf{R} \in \mathbb{R}^{d \times d}$ for intermixing outputs across memory cells~\citep{Greff:15}; further exposition is provided in Appendix~\ref{sec:apporgLSTM}. Empirical ablation results have confirmed the indispensable nature of each element within the memory cell structure~\citep{Greff:15}.

\subsection{sLSTM}
\label{sec:sLSTM}

To provide LSTMs with the capability to reconsider their storage choices, we incorporate exponential gates (red) along with techniques for normalization and stabilization. Specifically, we allow the input and forget gates to utilize exponential activation functions. Regarding normalization, we implement a normalizer state that accumulates the sum of the products between the input gate and every subsequent forget gate.

The scalar sLSTM forward pass is:

\begin{align}
\label{eq:slstmforward}
c_t \ &= \  f_t \ c_{t-1} \ + \ i_t \ z_t & & 
 &\text{cell state} \\
n_t \ &= \  f_t \ n_{t-1} \ + \ i_t  & & 
 &\text{normalizer state} \\
h_t \ &= \ o_t \ \tilde{h}_t \ ,
  &\tilde{h}_t \ &= \ c_t / n_t
&\text{hidden state} \\
z_t \ &= \ \varphi \left( \tilde{z}_t \right) \ , 
  &\tilde{z}_t \ &=  \ \mathbf{w}^\top_{z} \ \mathbf{x}_t \ + \
  r_{z}  h_{t-1} \ + \  b_{z}
  &\text{cell input} \\
i_t \ &= \ \!\exp \left( \tilde{i}_t  \right) \ , 
  &\tilde{i}_t \ &= \ \mathbf{w}^\top_{i} \ \mathbf{x}_t \ + \
  r_{i}  \ h_{t-1} \ + \  b_{i}
  &\text{input gate} \\
f_t \ &= \ \sigma \left(  \tilde{f}_t \right) \ \text{OR} \
\!\exp \left(  \tilde{f}_t \right) \ ,
  &\tilde{f}_t \ &= \ \mathbf{w}^\top_{f} \ \mathbf{x}_t  \ + \
  r_{f}  \ h_{t-1} \ + \  b_{f}
  &\text{forget gate} \\
o_t \ &= \ \sigma \left( \tilde{o}_t \right) \ , 
  &\tilde{o}_t  \ &= \ \mathbf{w}^\top_{o} \ \mathbf{x}_t \ + \
  r_{o}  \ h_{t-1} \ + \  b_{o}
  &\text{output gate} 
\end{align}

We transfer the original LSTM gating techniques, i.e., input- and/or hidden-dependent gating plus bias term, to the new architectures. Exponential activation functions can lead to large values that cause overflows. Therefore, we stabilize gates with an additional state \mbox{$m_t$~\citep{Milakov:18arxiv}}:

\begin{align}
\label{eq:slstmstabil}
m_t \ &= \ \max \left( \log ( f_t ) + m_{t-1} , \log ( i_t ) \right) &\text{stabilizer state} \\
i'_t \ &= \ \!\exp \left( \log \left ( i_t \right) - m_t \right) \ = \ \!\exp \left( \tilde{i}_t - m_t \right) \ \, 
  &\text{stabil. input gate} \\
f'_t \ &= \ \!\exp \left( \log \left( f_t \right) + m_{t-1} - m_t \right) \ \, 
  &\text{stabil. forget gate}
\end{align}

As demonstrated in Appendix~\ref{sec:appsLSTM}, substituting $f_t$ with $f'_t$ and $i_t$ with $i'_t$ during the forward computation leaves both the overall network output and the loss gradients with respect to the parameters unaffected.

\paragraph{New Memory Mixing.} Similar to the standard LSTM, sLSTM supports having several memory cells (refer to Appendix~\ref{sec:appsLSTM}). The presence of multiple cells facilitates memory mixing by means of recurrent links $\mathbf{R}_{\mathbf{z}}$, $\mathbf{R}_{\mathbf{i}}$, $\mathbf{R}_{\mathbf{f}}$, and $\mathbf{R}_{\mathbf{o}}$ running from the hidden state $\mathbf{h}$ to each of the memory cell input $\mathbf{z}$ and gate vectors $\mathbf{i}$, $\mathbf{f}$, and $\mathbf{o}$. A novel dimension in this mixing mechanism comes from applying exponential gating. With the updated sLSTM, multiple heads are permitted, where memory mixing occurs internally within each head but not between different heads. By combining the head structure in sLSTM with exponential gating, this approach realizes an innovative memory mixing strategy.

\subsection{mLSTM}
\label{sec:mLSTM}

In order to boost the storage capabilities of LSTM networks, we substitute the conventional scalar memory cell $c \in \mathbb{R}$ with a matrix-based cell $\mathbf{C} \in \mathbb{R}^{d \times d}$. This change shifts retrieval operations to matrix multiplications. At each time step $t$, the model is tasked with storing two vectors: a key $\mathbf{k}_t \in \mathbb{R}^d$ and a value $\mathbf{v}_t \in \mathbb{R}^d$, adopting the nomenclature from Transformers. At a subsequent time step $t + \tau$, we aim to recover the value vector $\mathbf{v}_t$ using a query $\mathbf{q}_{t+\tau} \in \mathbb{R}^d$. This framework corresponds to the paradigm of Bidirectional Associative Memories (BAMs) \citep{Kohonen:72,Anderson:72,Nakano:72,Anderson:77}.

The covariance update rule~\citep{Sejnowski:77,Dayan:91} for storing a key-value pair is

\begin{align}
\mathbf{C}_t \ &= \ \mathbf{C}_{t-1} \ + \ \mathbf{v}_{t} \ \mathbf{k}_t^\top \ .
\end{align}

We consider inputs that have undergone layer normalization before being projected onto the key and value spaces, ensuring they maintain zero mean. The covariance update mechanism has been shown to be optimal for maximizing the separability of retrieved binary representations, corresponding to an optimal signal-to-noise ratio~\citep{Dayan:91}. Restricting retrieval dynamics to pairwise terms and accepting the quadratic computational burden characteristic of attention-based methods can yield even higher levels of separability~\citep{Krotov:16,Krotov:17,Ramsauer:21}. Notably, the covariance update formulation aligns with the Fast Weight Programmer paradigm~\citep{Schmidhuber:92ncfastweights,Schlag:21}, which was subsequently modified to incorporate a fixed decay factor applied to $\mathbf{C}_{t-1}$ and a fixed learning rate applied to 
$\mathbf{v}_{t} \mathbf{k}_t^\top$~\citep{Ba:16}.
Drawing from this line of reasoning, we embed the covariance update procedure within the LSTM architecture, interpreting the forget gate as the decay parameter, the input gate as the learning rate, and the output gate as a scaling factor on the recovered vector.

For this matrix memory, the normalizer state is the weighted sum of key vectors, where each key vector is weighted by the input gate and all future forget gates. Again, the normalizer state keeps record of the strength of the gates. Since the dot product between query and normalizer state can be close to zero, we use the absolute value of this dot product and lower bound it by a threshold (typically 1.0) as done previously~\citep{Sun:23arxiv}. The mLSTM forward pass is:

\begin{align}
\label{eq:mlstm_recurrent_begin}
\mathbf{C}_t \ &= \  f_t \ \mathbf{C}_{t-1} \ + \ 
  i_t \ \mathbf{v}_t \ \mathbf{k}_t^\top & &  &\text{cell state} \\
\mathbf{n}_t \ &= \  f_t \ \mathbf{n}_{t-1} \ + \ 
  i_t \ \mathbf{k}_t & &  &\text{normalizer state} \\
\mathbf{h}_t  \ &= \ \mathbf{o}_t \ \odot \ \tilde{\mathbf{h}}_t
\ , \qquad \qquad 
\tilde{\mathbf{h}}_t \ = \   \mathbf{C}_t \mathbf{q}_t \ / \ 
  \max \left\{ |\mathbf{n}_t^\top \mathbf{q}_t|, 1 \right\} 
   &&&\text{hidden state} \\
\mathbf{q}_t \ &= \ \mathbf{W}_q \ \mathbf{x}_t \ + \ \mathbf{b}_q  & &  &\text{query input} \\
\mathbf{k}_t \ &= \ \frac{1}{\sqrt{d}} \mathbf{W}_k \ \mathbf{x}_t \ + \ \mathbf{b}_k  & &  &\text{key input} \\
\mathbf{v}_t \ &= \ \mathbf{W}_v \ \mathbf{x}_t \ + \ \mathbf{b}_v  & &  &\text{value input} \\
i_t \ &= \ \!\exp \!  \! \left( \tilde{i}_t  \right) \ , 
 \qquad \qquad \ \ \ \ \,
  \tilde{i}_t \ = \ \mathbf{w}^\top_{i} \ \mathbf{x}_t \ + \  b_{i} \ \ \ 
  &&&\text{input gate} \\
f_t \ &= \ \sigma \!  \left(  \tilde{f}_t \right) \ \text{OR} \ \!\exp \!  \! \left(  \tilde{f}_t \right) , \ \ \: \!
\tilde{f}_t \ = \ \mathbf{w}^\top_{f} \ \mathbf{x}_t  \ + \
  b_{f}
  &&& \text{forget gate} \\
\label{eq:mlstm_recurrent_end}
\mathbf{o}_t \ &= \ \sigma \left( \tilde{\mathbf{o}}_t \right) \ , \qquad \qquad \quad \ \ \ \,
  \tilde{\mathbf{o}}_t  \ = \ \mathbf{W}_{\mathbf{o}} \ \mathbf{x}_t \ + \
  \mathbf{b}_{\mathbf{o}}
  &&&\text{output gate} 
\end{align}

mLSTM, similar to standard LSTM architectures, supports several memory cells. In the mLSTM framework, having multiple heads or multiple cells yields an equivalent structure due to the absence of interactions between memory units. To ensure stable behavior of the exponential gating mechanisms within mLSTM, we employ stabilization strategies identical to those implemented in sLSTM (refer to Equation~\ref{eq:slstmstabil}). 
Given that mLSTM lacks any form of memory mixing, its recurrence relation can be rewritten in a manner that facilitates parallel computation. Additional information can be found in Appendix~\ref{sec:appmLSTM}.

\subsection{xLSTM Architecture}
\label{sec:xLSTMarch}

\paragraph{xLSTM Blocks.}
\label{sec:xLSTMblock}
The purpose of an xLSTM block is to non-linearly encode previous information within a high-dimensional representation, thereby enhancing the distinction between varying historical sequences or contextual backgrounds. Such differentiation among previous histories is essential for accurate prediction of the subsequent sequence element, for example, forecasting the next token. Our approach is grounded in Cover’s Theorem~\citep{Cover:65}, which posits that patterns transformed into a higher-dimensional space via nonlinear embedding become more readily separable through linear means compared to their original representations. We evaluate two forms of residual block configurations: (i) A residual block featuring a post up-projection, akin to those found in Transformer architectures. Here, a nonlinear summary is formed in the original space, followed by a linear projection into a higher-dimensional setting, after which a nonlinear activation is applied, and finally, a linear mapping returns the result to the original dimensionality; refer to the leftmost panel in Figure~\ref{fig:Backbones} and the third column of Figure~\ref{fig:xlstm_sketch}. For a comprehensive illustration, see Figure~\ref{fig:appsLSTM_detailed} in the appendix. (ii) A residual block employing a pre up-projection, drawing parallels to State Space Models, initiates by first linearly projecting into a high-dimensional space,

summarizes historical information in a non-linear manner within a high-dimensional latent space, subsequently projecting it back to the original dimensionality via a linear mapping. When an xLSTM block incorporates an sLSTM, we employ the post up-projection module. Conversely, if an xLSTM block utilizes an mLSTM, we opt for the pre up-projection mechanism, taking into account the increase in memory capacity present in the higher-dimensional representation. For a comprehensive illustration, refer to the leftmost panel of Figure~\ref{fig:Backbones}, the third column of Figure~\ref{fig:xlstm_sketch}, or Figure~\ref{fig:appsLSTM_detailed} included in the appendix.

\paragraph{xLSTM Architecture. }
\label{sec:xLSTMarchitecure}
The xLSTM model is formed by stacking its fundamental modules in a residual manner, following approaches outlined by~\citet{Srivastava:15,He:16}. The architecture adopts a standard residual backbone with pre-LayerNorm~\citep{Ba:16b}, which is a prevalent choice in current Large Language Model designs. Refer to the final column of Figure~\ref{fig:xlstm_sketch} for visualization.

\subsection{Memory and Speed Considerations}

In contrast to Transformers, xLSTM architectures exhibit linear computational requirements and maintain constant memory usage regardless of sequence length. The inherent compressive nature of xLSTM memory makes it particularly advantageous for use in industrial contexts and for deployment on edge devices.

Unlike conventional approaches, the mLSTM model leverages parameter-free memory but incurs significant computational costs arising from the need for a $d \times d$ matrix to store memory and perform updates. This design strikes a balance between the available memory capacity and the necessary computation. Importantly, these matrix operations are amenable to parallel processing on GPUs, thus contributing minimally to overall wall-clock runtime.

Although mLSTM lends itself to parallel execution in a fashion similar to techniques like FlashAttention~\citep{Dao:22,Dao:23} or GLA~\citep{Yang:23arxiv}, the sLSTM variant cannot be similarly parallelized because of its use of memory mixing via hidden-to-hidden interactions. To address this, we have engineered an efficient CUDA implementation with memory optimizations down to the GPU register level, resulting in execution times that are generally within a factor of two of mLSTM’s speed.

\section{Related Work}
\label{sec:related}

\paragraph{Linear Attention.} To address the issue of quadratic computational complexity with respect to context length in the Transformer architecture, numerous approaches have been introduced that render the attention mechanism linear in this dimension. The Synthesizer generates synthetic attention weights independently of token--token relationships~\citep{Tay:20arxiv}. Linformer employs a low-rank matrix representation to approximate self-attention in a linear fashion~\citep{Wang:20arxiv}. Linear Transformer transforms the attention mechanism itself into a linear operation~\citep{Katharopoulos:20}. Performer utilizes positive orthogonal random features to efficiently approximate the attention softmax linearly~\citep{Choromanski:21}. Additionally, alternatives to standard attention—such as Structured Global Convolution (SGConv)~\citep{Li:22} and the Hyena Hierarchy~\citep{Poli:23}—substitute attention operations with rapid long-range convolutions.

\paragraph{State Space Models.} In recent times, State Space Models (SSMs) have garnered significant attention due to their linear complexity with respect to input sequence length and their competitive results in comparison to Transformer-based architectures. Among the initial contributions in this area was the Structured State Space sequence model (S4)~\citep{Gu:21}, which was subsequently followed by the introduction of the Diagonal State Space (DSS) model~\citep{Gupta:22}, Gated State Space (GSS) models~\citep{Mehta:22}, the S5 model~\citep{Smith:22}, Bidirectional Gated SSM (BiGS)~\citep{Wang:22}, H3 model~\citep{Fu:23}, and, more recently, Mamba~\citep{Gu:24arxiv}.

\paragraph{Recurrent Neural Networks.} Recurrent Neural Networks (RNNs) have been proposed as alternatives to Transformers and attention mechanisms, largely because their computational cost scales linearly with sequence length. Notably, RNNs built using Deep Linear Recurrent Units (LRUs) have demonstrated strong performance on language modeling tasks~\citep{Orvieto:23, De:24arxiv}. Similarly, the Hierarchically Gated Linear RNN (HGRN)~\citep{Qin:23} and its successor HGRN2~\citep{Qin:24arxiv} have yielded favorable results. Among the RNN-based strategies, RWKV~\citep{Peng:23arxivshort,Peng:24arxivshort} is recognized for attaining results comparable to those of Transformer architectures in large-scale language modeling.

\paragraph{Gating.}
A foundational concept introduced in LSTM architectures is gating, a mechanism that has since been independently revisited and integrated across various contemporary methods. This concept forms an essential component in frameworks such as HGRN~\citep{Qin:23}, HGRN2~\citep{Qin:24arxiv}, Gated Linear Attention (GLA)~\citep{Yang:23arxiv}, Gated State Space (GSS) models~\citep{Mehta:22}, Bidirectional Gated SSM (BiGS)~\citep{Wang:22}, Moving Average Equipped Gated Attention (MEGA)~\citep{Ma:22}, RWKV~\citep{Peng:23arxivshort}, and Mamba~\citep{Gu:24arxiv}.

\paragraph{Covariance Update Rule.}
To improve the ability of mLSTM cells to store information, we introduce a matrix-based memory governed by a covariance update rule. Similar update-based strategies appear in approaches such as Fast Weight Programmers \citep{Schmidhuber:92ncfastweights,Schlag:21}, RWKV-5 and RWKV-6~\citep{Peng:24arxivshort}, Retention~\citep{Sun:23arxiv}, Linear Transformer~\citep{Katharopoulos:20}, as well as HGRN2~\citep{Qin:24arxiv}.

\paragraph{Most Related.} The models most akin to xLSTM from a conceptual standpoint include Retention~\citep{Sun:23arxiv}, RWKV~\citep{Peng:23arxivshort,Peng:24arxivshort}, and HGRN2~\citep{Qin:24arxiv}, as they incorporate mechanisms such as matrix memory and gating. Nevertheless, unlike the recently introduced sLSTM, these alternatives lack support for memory mixing. The presence of memory mixing facilitates the resolution of state tracking challenges, granting LSTMs a greater level of expressiveness compared to both State Space Models (SSMs) and Transformers~\citep{Merrill:24,Deletang:23}. Effective state tracking proves crucial for tasks such as code execution or maintaining continuity of entities throughout extended narratives.

\paragraph{Residually Stacking Architectures.} xLSTM architectures, consistent with nearly all leading-edge deep learning frameworks, are built using a residual stacking of core modules~\citep{Srivastava:15,He:16}.

This modular residual approach has underpinned the advances in deep convolutional neural networks~\citep{He:16} as well as in Transformers~\citep{Vaswani:17}. The Transformer architecture, in particular, is foundational to the success of Large Language Models (LLMs) such as GPT-3~\citep{Brown:20short}, ChatGPT~\citep{Schulman:22short}, GPT-4~\citep{Achiam:23gpt4short}, Megatron-LM~\citep{Shoeybi:19}, Gopher~\citep{Rae:21short}, ERNIE 3.0 Titan~\citep{Wang:21arxivshort}, GLaM~\citep{Du:21glamshort}, Chinese M6~\citep{Lin:21}, the multilingual AlexaTM 20B~\citep{Soltan:22}, OPT~\citep{Zhang:22opt}, Chinchilla~\citep{Hoffmann:22short}, BLOOM~\citep{Scao:22arxivshort}, GLM-130B~\citep{Zeng:22arxivshort}, LaMDA~\citep{Thoppilan:22short}, PaLM~\citep{Chowdhery:22short}, Llama~\citep{Touvron:23arxiv}, and Gemini~\citep{GeminiTeam:23arxiv,Reid:24arxiv}.

\section{Experiments}
\label{sec:Exp}

In this section, we present an empirical assessment of xLSTM, benchmarking its performance against existing approaches with an emphasis on language modeling applications. The particular strengths of xLSTM on artificial benchmarks are explored in Section~\ref{sec:ExpSyntethic}. 
Section~\ref{sec:ExpComparison} is devoted to evaluating the validation set perplexity achieved by multiple leading language modeling techniques, all trained on a 15B-token subset of SlimPajama~\citep{Soboleva:23}. Here, we additionally report results from ablation studies conducted on xLSTM within the same data regime. The scaling laws exhibited by each method are then analyzed following the frameworks of \citet{Kaplan:20} and \citet{Brown:20short}.

Proceeding to Section~\ref{sec:ExpLanguage}, 
we expand our investigation with a comprehensive language modeling study. In this part, xLSTM and the top-performing techniques identified in Section~\ref{sec:ExpComparison} are trained on 300B tokens from SlimPajama~\citep{Soboleva:23}, allowing for a deeper comparison. We begin by measuring each model’s capacity for handling longer context extrapolation, then proceed to evaluate them via validation perplexity and their results on downstream benchmarks~\citep{Sutawika:24}. Furthermore, we assess generalization by evaluating performance on the 571 textual domains of the PALOMA language benchmark dataset~\citep{Magnusson:23arxivshort}, and finally, we revisit the scaling analysis for all models, this time under a training regime utilizing 20 times more data than before.

\label{sec:xlstm_blocks_notation}
Throughout all experiments, we denote xLSTM[$a$:$b$] as representing a configuration with a ratio of $a/b$ between mLSTM-based and sLSTM-based xLSTM blocks. For instance, xLSTM[7:1] specifies a total of eight blocks, seven of which employ mLSTM while one utilizes sLSTM. When the overall number of blocks is 48, this corresponds to 42 mLSTM-based and 6 sLSTM-based blocks. Additionally, across all experiments, mLSTM blocks are paired with pre up-projection blocks, and sLSTM blocks are paired with post up-projection blocks.

\subsection{Synthetic Tasks and Long Range Arena}
\label{sec:ExpSyntethic}
We begin by evaluating xLSTM's novel exponential gating mechanism combined with memory mixing on a suite of formal language benchmarks~\citep{Deletang:23}. Subsequently, we investigate how well xLSTM's innovative matrix memory structure handles the Multi-Query Associative Recall task~\citep{Arora:23arxiv}. Lastly, we analyze xLSTM's capabilities on long-context processing using the Long Range Arena benchmarks~\citep{Tay:21}.

\paragraph{Evaluation of xLSTM's Exponential Gating with Memory Mixing.} We evaluate the newly introduced exponential gating with memory mixing in xLSTM, designed to address challenges in state tracking tasks~\citep{Merrill:24, Merrill:23}. Building on the formal language benchmarks from \citet{Deletang:23}, we extend them to permit multi-length training, facilitating length extrapolation capabilities. Details regarding the tasks, as well as comprehensive experimental results, can be found in Appendix~\ref{sec:appExpSynthetic-formalLang}. For these tasks, xLSTM is assessed against a range of baselines, including Transformers, State Space Models, and Recurrent Neural Networks. The evaluation metric focuses on token-wise accuracy pertinent to each task, normalized so that 0 corresponds to random prediction and 1 to perfect performance. Specifically, we test 2-block variants of the following approaches: xLSTM[0:1] (using only sLSTM), xLSTM[1:0] (using only mLSTM), xLSTM[1:1], Llama, Mamba, RWKV, Retention, Hyena, LSTM, and LSTM within Transformer blocks (LSTM (Block)). A summary of the experimental outcomes is presented in Figure~\ref{fig:formal-main}. Notably, models lacking memory mixing—such as Transformers or State Space Models—fail to resolve tasks requiring state tracking, including regular grammar challenges like the parity problem. These observations align with prior work demonstrating that Transformers and State Space Models are intrinsically less expressive than RNNs~\citep{Merrill:24, Merrill:23, Deletang:23}.

\paragraph{Test of xLSTM's Memory Capacities on Associative Recall Tasks.} This study evaluates the capacity of xLSTM’s matrix memory using the Multi-Query Associative Recall task~\citep{Arora:23arxiv}. For each test sequence, a set of key-value pairs is selected at random from an extensive vocabulary, with the requirement that these associations be remembered for later retrieval. To intensify the challenge compared to the original formulation, we increase both the number of key-value associations (up to 256) and the length of the input context (up to 2048). These adjustments enable a more extensive assessment of different models’ memory capabilities. Our comparison includes 2-block variants of Llama, Mamba, RWKV-5, RWKV-6, xLSTM[1:1], and xLSTM[1:0]. Model performance is measured by the proportion of correctly recalled pairs. Owing to their exponential memory scaling with respect to the embedding dimension~\citep{Ramsauer:21}, Transformer models like Llama serve as a benchmark for this task. The main findings are summarized in Figure~\ref{fig:mqar-main}. 
Among non-Transformer approaches, xLSTM[1:1] consistently achieves the highest recall accuracy, regardless of model size. Notably, the inclusion of the sLSTM block does not impede overall memory performance—in fact, it enhances it, especially apparent in the most challenging setting with 256 key-value pairs. Supplementary results can be found in Appendix~\ref{sec:appExpSynthetic-mqar}, where further analyses suggest that xLSTM’s improved memory structure enables effective generalization to input lengths surpassing those encountered during model training.

\paragraph{Test of xLSTM's Long Context Capabilities on Long Range Arena.} In order to evaluate xLSTM's ability to process lengthy input sequences and maintain extended contextual information, we benchmark various approaches using the Long Range Arena~\citep{Tay:21}. Across all evaluated tasks, xLSTM exhibits robust and reliable results, indicating that its architectural design effectively manages the diverse challenges presented by long context scenarios.

For more details, see Appendix~\ref{sec:appExpSynthetic-lra}.

\subsection{Method Comparison and Ablation Study}
\label{sec:ExpComparison}

In order to investigate the central issue posed in this work—specifically, the capabilities of our novel LSTM variants when applied to large-scale language modeling—we conduct experiments where xLSTMs, Transformers, State Space Models, and several other approaches are all trained with 15B tokens sourced from SlimPajama, employing an identical auto-regressive framework. Evaluation is carried out using the validation dataset, and comprehensive ablation studies are conducted on xLSTMs to further analyze their properties.

\paragraph{Comparison of xLSTM with Alternative Approaches.} To conduct a fair evaluation, all models are trained using 15B tokens sourced from SlimPajama~\citep{Soboleva:23}, and their validation set perplexities are subsequently assessed. The following approaches are included in the comparison: xLSTM (the introduced approach), GPT-3 (Transformer)~\citep{Brown:20short}, Llama (Transformer)~\citep{Touvron:23arxiv}, H3 (SSM)~\citep{Fu:23}, Mamba (SSM)~\citep{Gu:24arxiv}, RWKV-4, RWKV-5, RWKV-6 (RNNs)~\citep{Peng:23arxivshort, Peng:24arxivshort}, GLA (linear Transformer)~\citep{Yang:23arxiv}, HGRN, HGRN2 (RNNs)~\citep{Qin:23, Qin:24arxiv}, RetNet (linear Transformer)~\citep{Sun:23arxiv}, Hyena (linear Transformer)~\citep{Poli:23}, as well as xLSTM[1:0] and xLSTM[7:1] (see Section~\ref{sec:xlstm_blocks_notation} for notation). Training utilized mixed precision across models; however, for RWKV-5, RWKV-6, GLA, and HGRN2, PyTorch's automated mixed precision was not adopted (see Appendix Section~\ref{sec:appGenTrainProc} for clarification).

The approaches can be grouped as follows: (a) Transformers, (b) State Space Models (SSMs), and (c) Recurrent Neural Networks (RNNs), alongside linear Transformer variants. The latter are distinguished by their linearized alternatives to the conventional Transformer attention mechanism. Each model is designed to be comparable in scale to a 350M-parameter GPT-3 instance, employing a 1024-dimensional embedding and comprising 24 residual blocks. Notably, only the GPT-3 implementation makes use of shared token and output embedding weights, reducing its parameter count relative to the others.

Results presented in Table~\ref{tab:spaj15b_model_results} clearly indicate that xLSTM achieves the lowest validation perplexity, surpassing all baseline methods. The Appendix~\ref{sec:appExpComparison} provides more granular information. Additionally, Figure~\ref{fig:temp_sclaw15B} illustrates the scaling properties for these experiments, demonstrating that xLSTM continues to deliver competitive performance as model size increases.

\paragraph{Ablation Studies.}
As shown in Table~\ref{tab:spaj15b_model_results} and Figure~\ref{fig:temp_sclaw15B}, the xLSTM model demonstrates superior language modeling capabilities when trained with 15B tokens from SlimPajama. To systematically investigate the contributions of architectural modifications from LSTM to xLSTM, we incrementally transform a standard LSTM architecture into that of xLSTM. The procedure begins by embedding LSTM layers into a pre-LayerNorm residual backbone. Next, we enhance the model with a post up-projection block. The final modifications involve incorporating exponential gating and matrix memory mechanisms.

The results are shown in Table~\ref{tab:ablstudies} (top). 

The ablation experiments demonstrate that the notable performance gains arise from the combination of exponential gating and matrix memory. Furthermore, given the critical role of gating in both RNNs and State Space Models, we systematically remove or modify various gating strategies. As reported in Table~\ref{tab:ablstudies} (bottom), we observe that making each gate learnable and responsive to input produces cumulative improvements. Further detailed analyses of the backbone's individual elements can be found in Appendix~\ref{sec:appAblDetails}.

\subsection{xLSTM as Large Language Model}
\label{sec:ExpLanguage}

To conclude this work, we conduct extensive language modeling experiments at scale to evaluate the viability of xLSTM as a large language model (LLM). Specifically, we enlarge the training corpus, utilizing 300B tokens drawn from SlimPajama. This quantity aligns with that employed by, for instance, Mamba~\citep{Gu:24arxiv} and Griffin~\citep{De:24arxiv}.

For comparative analysis, we benchmark xLSTM against RWKV-4, Llama, and Mamba, with each chosen to represent distinct architectural classes as outlined in Section~\ref{sec:ExpComparison}. RWKV-4 serves as the RNN baseline since proper training precision for RWKV-5, RWKV-6, and HGRN2 became established only after the commencement of our 300B token experiments (refer to Appendix~\ref{sec:appGenTrainProc}).

Models are trained at various scales (125M, 350M, 760M, and 1.3B parameters). We systematically evaluate every configuration for their capacity to extrapolate to longer input sequences, measure their effectiveness on the validation set, and further probe their performance across downstream tasks. In addition, we subject the models to language modeling evaluation on 571 domains from the PALOMA benchmark and analyze their adherence to scaling laws.

\paragraph{Sequence Length Extrapolation.}
\label{sec:spaj300B_ctx_extrapolation}
To begin, we assess how well large-scale (1.3B parameter) xLSTM, RWKV-4, Llama, and Mamba models extrapolate to longer sequence lengths. These models are each trained on a context window of 2048 tokens and subsequently evaluated using increasingly longer context lengths, reaching up to 16384 tokens. The findings, presented in Figure~\ref{fig:spaj300B_seq_extrapolation}, reveal that xLSTM consistently yields lower perplexities than the other architectures as context length increases.

\paragraph{Validation Perplexity and Downstream Tasks.}
\label{sec:spaj_downstream_eval}
Next, across all examined parameter scales, we measure how xLSTM, RWKV-4, Llama, and Mamba models perform on the SlimPajama validation dataset for next-token prediction and evaluate their abilities on downstream benchmarks focusing on common sense reasoning.

The third column in Table~\ref{tab:lmeval} shows the validation set perplexities across various approaches. xLSTM[1:0] and xLSTM[7:1] consistently achieve the lowest perplexity on the validation set for every model size evaluated. The remaining columns in Table~\ref{tab:lmeval} summarize downstream task performance. In nearly all tasks and for all model capacities, xLSTM outperforms competing methods; the only exception occurs in certain instances on the ARC benchmark, where Mamba surpasses the other approaches. Additional discussion is presented in Appendix~\ref{sec:appExpLanguage}.

\paragraph{Performance on PALOMA Language Tasks.} As a third step, we evaluate the capabilities of xLSTM, RWKV-4, Llama, and Mamba models in next token prediction across all model capacities using the PALOMA language tasks~\citep{Magnusson:23arxivshort}. Evaluation is performed via next token perplexity on 571 unique text domains, spanning sources such as nytimes.com as well as Reddit communities like r/depression.

Table~\ref{tab:ppleval} presents the resulting perplexity figures, distinguishing between broader language modeling tasks (first seven columns) and more specialized domain benchmarks (last five columns). Notably, xLSTM[1:0] outperforms xLSTM[7:1] within this language evaluation.

Across these 571 text domains, xLSTM[1:0] achieves lower perplexity than Mamba in 568 domains (99.5\%), surpasses Llama in 486 domains (85.1\%), and outperforms RWKV-4 in 570 domains (99.8\%). Further details can be found in Appendix~\ref{sec:appExpLanguage}.

\paragraph{Scaling Laws.} Fourth, we investigate the presence of power-law scaling, which provides a basis for predicting how performance extends as model size increases~\citep{Kaplan:20,Brown:20short}. Figure~\ref{fig:temp_sclaw300B} illustrates these scaling dynamics. While all models exhibit broadly similar scaling trends, they do so at distinct baselines. RWKV-4 demonstrates the weakest performance, with Llama and Mamba following. xLSTM outperforms Mamba, maintaining a relative improvement comparable to the gap observed between Mamba and Llama. The observed scaling trends suggest that xLSTM is poised to retain its performance edge over Transformer and State-Space architectures as model size grows.

\textbf{Generation Times and Maximal Throughput.} In our evaluation, depicted in Figure~\ref{fig:inference_time_generation}, we report on text generation latencies, and in Figure~\ref{fig:inference_time_decoding_throughput} (left), we examine peak throughput for the xLSTM variants at the 1.3B parameter scale. Our comparison set includes equivalently sized implementations of Mamba, Llama, and RWKV from HuggingFace, utilizing a static key-value cache for the Llama configuration. It should be noted that, during these experiments, compilation techniques such as full cache compilation for the Transformer model and application of \texttt{torch.compile} to the Mamba model were not functioning. All text generation benchmarks employ a batch size of 1 with a pre-fill of 16 tokens, a scenario that particularly benefits Transformer-based architectures. As illustrated in Figure~\ref{fig:inference_time_generation}, both xLSTM and the recurrent models (Mamba and RWKV-4) display linear time complexity in generation, contrasting the quadratic time complexity evidenced by Llama. Throughput measurements further explore a range of batch sizes and pre-fill values for Llama. From Figure~\ref{fig:inference_time_decoding_throughput} (right), it is evident that xLSTM accommodates substantially larger batch sizes compared to Llama, owing to its fixed memory demands, and thereby realizes superior throughput.

\section{Limitations}

(i) Unlike mLSTM, the memory mixing mechanism in sLSTM inhibits the use of parallelizable computation, which prevents rapid parallel implementations. Despite this challenge, we have created an efficient CUDA kernel for sLSTM that achieves speeds within a factor of two compared to our parallel mLSTM implementation. (ii) Presently, the CUDA kernels for mLSTM are not fine-tuned for performance, resulting in runtimes approximately four times longer than those of FlashAttention or the scan operation deployed in Mamba. Future improvements in CUDA kernel design—drawing inspiration from FlashAttention—could yield faster results. (iii) The matrix memory architecture in mLSTM involves handling $d \times d$ matrices, incurring significant computational complexity. However, since memory update and access eschew trainable parameters and leverage standard, easily parallelizable matrix operations, the extra time incurred in practice is marginal. (iv) Special care must be exercised when initializing the forget gates. (v) Because the matrix memory does not scale with sequence length, increasing context length may exceed the memory capacity for extended sequences. Nevertheless, we have not observed this to be a bottleneck for sequence contexts as long as 16k tokens; further discussion can be found in Section~\ref{sec:spaj300B_ctx_extrapolation}. (vi) Owing to the substantial compute requirements of large-scale language model training, neither the architecture nor the hyperparameters have been comprehensively optimized, particularly for larger xLSTM variants. We expect that realizing the full capacity of xLSTM will necessitate thorough and deliberate optimization.

\section{Conclusion}
Our investigation provides a partial response to a straightforward inquiry: What can be achieved in language modeling by scaling LSTM architectures to the billion-parameter range? The results thus far reveal that such scaling yields performance on par with prevailing architectures, including Transformers and State Space Models. Through the introduction of exponential gating, memory mixing, and a novel memory organization, we have extended LSTM to the xLSTM variant. Experimental findings show that xLSTM demonstrates competitive, and often superior, results compared to leading methodologies such as Transformers and State Space Models in language modeling benchmarks. Observed scaling trends suggest that as xLSTM models grow larger, they may emerge as strong alternatives to Transformer-based Large Language Models. Beyond language modeling, xLSTM exhibits substantial promise for applications in other domains, including Reinforcement Learning, Time Series Forecasting, and the simulation of physical phenomena.

\end{document}